% Options for packages loaded elsewhere
\PassOptionsToPackage{unicode}{hyperref}
\PassOptionsToPackage{hyphens}{url}
%
\documentclass[
  ignorenonframetext,
]{beamer}
\usepackage{pgfpages}
\setbeamertemplate{caption}[numbered]
\setbeamertemplate{caption label separator}{: }
\setbeamercolor{caption name}{fg=normal text.fg}
\beamertemplatenavigationsymbolsempty
% Prevent slide breaks in the middle of a paragraph
\widowpenalties 1 10000
\raggedbottom
\setbeamertemplate{part page}{
  \centering
  \begin{beamercolorbox}[sep=16pt,center]{part title}
    \usebeamerfont{part title}\insertpart\par
  \end{beamercolorbox}
}
\setbeamertemplate{section page}{
  \centering
  \begin{beamercolorbox}[sep=12pt,center]{part title}
    \usebeamerfont{section title}\insertsection\par
  \end{beamercolorbox}
}
\setbeamertemplate{subsection page}{
  \centering
  \begin{beamercolorbox}[sep=8pt,center]{part title}
    \usebeamerfont{subsection title}\insertsubsection\par
  \end{beamercolorbox}
}
\AtBeginPart{
  \frame{\partpage}
}
\AtBeginSection{
  \ifbibliography
  \else
    \frame{\sectionpage}
  \fi
}
\AtBeginSubsection{
  \frame{\subsectionpage}
}
\usepackage{amsmath,amssymb}
\usepackage{lmodern}
\usepackage{iftex}
\ifPDFTeX
  \usepackage[T1]{fontenc}
  \usepackage[utf8]{inputenc}
  \usepackage{textcomp} % provide euro and other symbols
\else % if luatex or xetex
  \usepackage{unicode-math}
  \defaultfontfeatures{Scale=MatchLowercase}
  \defaultfontfeatures[\rmfamily]{Ligatures=TeX,Scale=1}
\fi
\usetheme[]{metropolis}
% Use upquote if available, for straight quotes in verbatim environments
\IfFileExists{upquote.sty}{\usepackage{upquote}}{}
\IfFileExists{microtype.sty}{% use microtype if available
  \usepackage[]{microtype}
  \UseMicrotypeSet[protrusion]{basicmath} % disable protrusion for tt fonts
}{}
\makeatletter
\@ifundefined{KOMAClassName}{% if non-KOMA class
  \IfFileExists{parskip.sty}{%
    \usepackage{parskip}
  }{% else
    \setlength{\parindent}{0pt}
    \setlength{\parskip}{6pt plus 2pt minus 1pt}}
}{% if KOMA class
  \KOMAoptions{parskip=half}}
\makeatother
\usepackage{xcolor}
\newif\ifbibliography
\usepackage{longtable,booktabs,array}
\usepackage{calc} % for calculating minipage widths
\usepackage{caption}
% Make caption package work with longtable
\makeatletter
\def\fnum@table{\tablename~\thetable}
\makeatother
\setlength{\emergencystretch}{3em} % prevent overfull lines
\providecommand{\tightlist}{%
  \setlength{\itemsep}{0pt}\setlength{\parskip}{0pt}}
\setcounter{secnumdepth}{5}
\newlength{\cslhangindent}
\setlength{\cslhangindent}{1.5em}
\newlength{\csllabelwidth}
\setlength{\csllabelwidth}{3em}
\newlength{\cslentryspacingunit} % times entry-spacing
\setlength{\cslentryspacingunit}{\parskip}
\newenvironment{CSLReferences}[2] % #1 hanging-ident, #2 entry spacing
 {% don't indent paragraphs
  \setlength{\parindent}{0pt}
  % turn on hanging indent if param 1 is 1
  \ifodd #1
  \let\oldpar\par
  \def\par{\hangindent=\cslhangindent\oldpar}
  \fi
  % set entry spacing
  \setlength{\parskip}{#2\cslentryspacingunit}
 }%
 {}
\usepackage{calc}
\newcommand{\CSLBlock}[1]{#1\hfill\break}
\newcommand{\CSLLeftMargin}[1]{\parbox[t]{\csllabelwidth}{#1}}
\newcommand{\CSLRightInline}[1]{\parbox[t]{\linewidth - \csllabelwidth}{#1}\break}
\newcommand{\CSLIndent}[1]{\hspace{\cslhangindent}#1}
\ifLuaTeX
  \usepackage{selnolig}  % disable illegal ligatures
\fi
\IfFileExists{bookmark.sty}{\usepackage{bookmark}}{\usepackage{hyperref}}
\IfFileExists{xurl.sty}{\usepackage{xurl}}{} % add URL line breaks if available
\urlstyle{same} % disable monospaced font for URLs
\hypersetup{
  pdftitle={Predicting Droughts in the Amazon Basin based on Global Sea Surface Temperatures},
  pdfauthor={Dario LepkeDr.~Fabian Scheipl (Ludwig-Maximilians-University Munich)Dr.~Niklas Boers (Potsdam Institute for Climate Impact Research)},
  hidelinks,
  pdfcreator={LaTeX via pandoc}}

\title{Predicting Droughts in the Amazon Basin based on Global Sea
Surface Temperatures}
\author{Dario Lepke\newline Dr.~Fabian Scheipl
(Ludwig-Maximilians-University Munich)\newline Dr.~Niklas Boers (Potsdam
Institute for Climate Impact Research)}
\date{September 08, 2022}

\begin{document}
\frame{\titlepage}

\begin{frame}[allowframebreaks]
  \tableofcontents[hideallsubsections]
\end{frame}
\hypertarget{introduction}{%
\section{Introduction}\label{introduction}}

\begin{frame}{Motivation}
\protect\hypertarget{motivation}{}
\begin{itemize}
\item
  The Amazon basin is a key hotspot of biodiversity, carbon storage and
  moisture recycling
\item
  Hydrological extremes affect ecosystem and populations tremendously
\item
  Droughts in the Amazon rainforest can have severe biomass carbon
  impact
\item
  Severe Amazon drought in 2010 had total biomass carbon impact of 2.2
  PgC, affected area \(3 mio km^2\)
\end{itemize}
\end{frame}

\begin{frame}{Related work}
\protect\hypertarget{related-work}{}
\begin{itemize}
\tightlist
\item
  Ciemer et al. (2020) established an early warning indicator for water
  deficits in the central Amazon basin (CAB)
\end{itemize}

\begin{figure}

{\centering \includegraphics[width=0.5\linewidth]{../figures/cross-degree} 

}

\caption{Cross degree between sea surface temperature and continental rainfall anomalies. For each grid cell of sea surface temperature in the Atlantic and Pacific, the cross degree towards rainfall in the Central Amazon Basin (blue box) is shown, for a positive correlations and b negative correlations. Darker shading indicates a larger cross degree, implying a larger number of links and thus significant correlations with rainfall at more grid points in the Central Amazon Basin. Red areas outline coherent oceanic regions with a the 20\% highest cross degrees for positive correlations, found in the Southern Pacific Ocean (SPO) and Southern Tropical Atlantic Ocean (STAO), and b the 20\% highest cross degrees for negative correlations, found in the Central Pacific Ocean (CPO) and Northern Tropical Atlantic Ocean (NTAO) (Ciemer et al. (2020))}\label{fig:cross}
\end{figure}

\begin{itemize}
\tightlist
\item
  They investigate the correlation over time of NTAO and STAO and its
  relationship to droughts in the CAB
\end{itemize}
\end{frame}

\begin{frame}{Early warning signal}
\protect\hypertarget{early-warning-signal}{}
\begin{figure}

{\centering \includegraphics[width=0.75\linewidth]{../figures/early-warning-signals} 

}

\caption{Early-warning signal for droughts in the central Amazon basin. We compare the time evolution of the average cross-correlation of the Northern Tropical Atlantic Ocean (NTAO) and Southern Tropical Atlantic Ocean (STAO), given by the blue curve, with the standardized precipitation index (SPI, orange) of the central Amazon basin. Orange dips indicate a negative SPI with a threshold for severely dry periods (SPI -1, dotted red line). We expect a drought event within the following one and a half years whenever the average cross-correlation between NTAO and STAO SST anomalies falls below an empirically found threshold of -0.06. Green circles indicate a matching forecast based on the Atlantic SST correlation structure, with one false alarm in 2002 indicated by a grey circle, where the threshold is crossed, but no drought took place in the direct aftermath (see Discussion). The temporal evolution of the average cross-correlation shown here is smoothed using a Chebyshev type-I low-pass filter with a cutoff at 24 months (Ciemer et al. (2020)).}\label{fig:early}
\end{figure}
\end{frame}

\begin{frame}{Our Approach}
\protect\hypertarget{our-approach}{}
\begin{itemize}
\tightlist
\item
  Inspect spatial and temporal characteristics in raw data
\item
  Directly predict rain from SST
\item
  Use lasso and fused lasso
\item
  model evaluation with cross validation for time series
\end{itemize}
\end{frame}

\hypertarget{explorative-analysis}{%
\section{Explorative analysis}\label{explorative-analysis}}

\begin{frame}{The Data}
\protect\hypertarget{the-data}{}
\begin{itemize}
\tightlist
\item
  Rain data from CHIRPS (Climate Hazards Group InfraRed Precipitation
  woth Station data)
\item
  CHIRPS contains in-situ and satellite data
\item
  SST data from ERSST (Extended Reconstructed Sea Surface Temperature)
\item
  ERSST is reanalysis of observation data (made by ships and buoys for
  example), missing data filled by interpolation techniques
\item
  These are the same data sets as in Ciemer et al. (2020)
\end{itemize}
\end{frame}

\begin{frame}{Explorative Analysis Rain}
\protect\hypertarget{explorative-analysis-rain}{}
\begin{figure}

{\centering \includegraphics[width=0.5\linewidth]{ma-presentation_files/figure-beamer/CAB-1} 

}

\caption{Location of the area under study. The central amazon basin (CAB) spanning across 0,-10 latitude and -70,-55 longitude}\label{fig:CAB}
\end{figure}
\end{frame}

\begin{frame}{Precipitation, Mean and SD}
\protect\hypertarget{precipitation-mean-and-sd}{}
\begin{figure}

{\centering \includegraphics[width=0.5\linewidth]{ma-presentation_files/figure-beamer/mean-sd-precip-1} 

}

\caption{Mean and standard deviation at each location. The standard deviation was computed over the whole time period. The white line on the scale at the side of the plots indicates the mean of the respective quantity}\label{fig:mean-sd-precip}
\end{figure}
\end{frame}

\begin{frame}{Precipitation Glyph Plots}
\protect\hypertarget{precipitation-glyph-plots}{}
\begin{figure}

{\centering \includegraphics[width=0.75\linewidth]{ma-presentation_files/figure-beamer/glyph-scale-1} 

}

\caption{Glyph map of de-seasonalised and smoothed precipitation. The time series are scaled locally, ranges are not the same in all cells. The different ranges are given in color shades, where lighter shading indicates a larger range and darker shades smaller ranges.}\label{fig:glyph-scale}
\end{figure}
\end{frame}

\begin{frame}{Explorative analysis SST}
\protect\hypertarget{explorative-analysis-sst}{}
\begin{figure}

{\centering \includegraphics[width=0.5\linewidth]{ma-presentation_files/figure-beamer/mean-and-sd-sst-og-1} 

}

\caption{Mean and SD on the global map.The color scales show the mean for the shown variable as white.}\label{fig:mean-and-sd-sst-og}
\end{figure}
\end{frame}

\hypertarget{correlation-analysis}{%
\section{Correlation analysis}\label{correlation-analysis}}

\begin{frame}{Correlation analysis}
\protect\hypertarget{correlation-analysis-1}{}
\begin{itemize}
\tightlist
\item
  Here we get an overview over the general correlation structure of the
  complete data
\item
  We show correlation data for the original as well as the seasonally
  adjusted data
\item
  The seasonal component was removed by using the stl algorithm that
  separates the time series into
\end{itemize}

\[ Monthly \textit{ } Data = Seasonal + Trend + Remainder\] - Two time
series can appear correlated but after removing the seasonal component
the correlation vanishes
\end{frame}

\begin{frame}{Correlation plot original SST}
\protect\hypertarget{correlation-plot-original-sst}{}
\begin{figure}

{\centering \includegraphics[width=0.65\linewidth,height=0.65\textheight]{../pres-plots/og-data/corr-0} 

}

\caption{Correlation plot between SST and mean precipitation in the CAB for timelag 0}\label{fig:corr-0}
\end{figure}
\end{frame}

\begin{frame}{Correlation plot original SST}
\protect\hypertarget{correlation-plot-original-sst-1}{}
\begin{figure}

{\centering \includegraphics[width=0.65\linewidth,height=0.65\textheight]{../pres-plots/og-data/q-corr-0} 

}

\caption{Correlation plot between SST and mean precipitation in the CAB for timelag 0}\label{fig:corr-0q}
\end{figure}
\end{frame}

\begin{frame}{Correlation plot de-seasonalized SST}
\protect\hypertarget{correlation-plot-de-seasonalized-sst}{}
\begin{figure}

{\centering \includegraphics[width=0.65\linewidth,height=0.65\textheight]{../pres-plots/des-data/corr-0} 

}

\caption{Correlation plot between de-seasonalized SST and de-seasonalized mean precipitation in the CAB for timelag 0}\label{fig:des-corr-0}
\end{figure}
\end{frame}

\begin{frame}{Correlation plot de-seasonalized SST}
\protect\hypertarget{correlation-plot-de-seasonalized-sst-1}{}
\begin{figure}

{\centering \includegraphics[width=0.65\linewidth,height=0.65\textheight]{../pres-plots/des-data/q-corr-0} 

}

\caption{Correlation plot between de-seasonalized SST and de-seasonalized mean precipitation in the CAB for timelag 0}\label{fig:des-corr-0q}
\end{figure}
\end{frame}

\hypertarget{clustering}{%
\section{Clustering}\label{clustering}}

\begin{frame}{Motivation/ Overview}
\protect\hypertarget{motivation-overview}{}
\begin{itemize}
\item
  explorative analysis has shown spatial and temporal differences in the
  precipitation data
\item
  we explored this further using k-means clustering
\item
  steps: find optimal k via pca and gap statistic
\item
  apply k-means to original precipitation data
\item
  we compared k-means and k-medoid with and without PCA via the gap
  statistic
\item
  here show only k-means with PCA as it gave best results
\item
  applying the regression models to separate clusters might improve
  predictions
\item
  Using 3 principal components and 5 cluster centers with k-means gave
  best results on gap statistic
\item
  we want to explore this further using a PCA and k-means clustering
\item
  find optimal number of k with the gap statistic
\item
  applying the regression models to separate clusters might improve
  predictions
\end{itemize}
\end{frame}

\begin{frame}{\(k\)-means}
\protect\hypertarget{k-means}{}
\begin{itemize}
\tightlist
\item
  Our objective is to find \(k\) internally homogeneous and externally
  heterogeneous clusters
\item
  Similarity is measured by the euclidean distance
\end{itemize}

\begin{equation} 
d(x_i,x_{i'}) = \sum_{j=1}^p(x_{ij}-x_{i'j})^2=||x_i-x_{i'}||^2
(\#eq:eucl-dist)
\end{equation}

\begin{itemize}
\tightlist
\item
  And we want to minimize the sum of distances inside all clusters,
  given by:
\end{itemize}

\begin{equation} 
W(C) = \frac{1}{2} \sum_{k=1}^{K} \sum_{C(i)=k} \sum_{C(i')=k} ||x_i-x_{i'}||^2 
= \sum_{k=1}^K N_k \sum_{C(i)=k} ||x_i-\bar{x}_k ||^2
(\#eq:wss)
\end{equation}

where \(\bar{x} = (\bar{x}_{1k},...,\bar{x}_{pk})\) stands for the mean
vectors of the \(k\)th cluster and \(N_k = \sum_{i=1}^N I(C(i)=k)\).
\end{frame}

\begin{frame}{gap statistic}
\protect\hypertarget{gap-statistic}{}
\begin{itemize}
\tightlist
\item
  number of clusters has to be defined beforehand
\item
  we decided on the optimal number of \(k\) using the gap statistic
\item
  Let \(W_k\) be \(W(C)\) for fix \(k\)
\item
  We compare \(W_k\) from the precipitation data with average \(W^*_k\)
  from \(B\) Monte Carlo sampled data sets
\end{itemize}

\begin{equation}
Gap(k) = E\{log(W*_k) \} - log(W_k).
(\#eq:gap)
\end{equation}

\begin{itemize}
\tightlist
\item
  We choose \(k\) as smallest k such that
\end{itemize}

\begin{equation}
Gap(k) \geq Gap(k+1) - s_{k+1}
\end{equation}

\begin{itemize}
\tightlist
\item
  \(s_{k+1}\) is \(sd_k\sqrt{1+1/B}\), and sd the standard deviation of
  log(W*\_k)
\end{itemize}
\end{frame}

\begin{frame}{PCA}
\protect\hypertarget{pca}{}
\begin{itemize}
\tightlist
\item
  Before running k-means we center the precipitation data and apply a
  PCA to reduce the large number of correlated variables to a few
\item
  The new variables are linear combinations of the original variables
\item
  Here: Each variable is a month of precipitation data in the CAB
\end{itemize}
\end{frame}

\begin{frame}{Scree Plot, PCA after centering}
\protect\hypertarget{scree-plot-pca-after-centering}{}
\begin{center}\includegraphics[width=0.75\linewidth]{ma-presentation_files/figure-beamer/unnamed-chunk-10-1} \end{center}
\end{frame}

\begin{frame}{Screeplot}
\protect\hypertarget{screeplot}{}
\begin{itemize}
\tightlist
\item
  The ``elbow'' be observe in the screeplot suggest 3 or 4 principal
  components
\item
  The first 3 and 4 first PC explain 67.77 and 70.79 of the variance
  respectively.
\item
  We compare the gap statistic results for 3 and 4 PC
\end{itemize}
\end{frame}

\begin{frame}{Gap statistic results}
\protect\hypertarget{gap-statistic-results}{}
\begin{center}\includegraphics[width=0.75\linewidth]{ma-presentation_files/figure-beamer/unnamed-chunk-11-1} \end{center}

\begin{itemize}
\tightlist
\item
  The k-means gap statistic on the first 3 PC proposes 5 clusters - For
  4 PC, 13 clusters are chosen
\item
  We chose 5 clusters since the result on 3 PC appears to be clearer and
  5 clusters are more applicable than fitting the model evaluation on 13
  clusters
\end{itemize}
\end{frame}

\begin{frame}{Clustering results}
\protect\hypertarget{clustering-results}{}
\begin{figure}

{\centering \includegraphics[width=0.5\linewidth]{ma-presentation_files/figure-beamer/cluster-map-1} 

}

\caption{Spatial distribution of the found clusters in the CAB. We applied a centered PCA on the data and used 3 principal components before applying the k-means algorithm}\label{fig:cluster-map}
\end{figure}

\begin{itemize}
\tightlist
\item
  We find 5 clusters of different sizes
\item
  The found clusters are almost completely spatially coherent although
  we did not include any spatial dependencies in the clustering
\item
  Small exception is the ``island'' of cluster 1 (orange) inside cluster
  4 (blue) and on the edge on cluster 3 (green)
\item
  Usefulness of clustering can only be determined after model fitting on
  each cluster
\end{itemize}
\end{frame}

\hypertarget{the-lasso}{%
\section{The lasso}\label{the-lasso}}

\begin{frame}{Definition of the lasso}
\protect\hypertarget{definition-of-the-lasso}{}
\begin{itemize}
\tightlist
\item
  We now consider the lasso regression problem
\end{itemize}

\begin{equation}
\min_{\beta_0, \beta} \frac{1}{N} \sum_{i=1}^N w_il(y_i,\beta_0 + \beta^Tx_i) + \lambda [(1-\alpha)||\beta||_2^2/2 + \alpha||\beta||_1] (\#eq:glmnet)
\end{equation}

\begin{itemize}
\tightlist
\item
  In our setting n \textless\textless{} p, so lasso is natural choice
\item
  The problem is solved using coordinate descent
\item
  Due to the time dependencies in our data normal Cross Validation may
  be unjustified
\end{itemize}
\end{frame}

\begin{frame}{Model evaluation}
\protect\hypertarget{model-evaluation}{}
\begin{itemize}
\tightlist
\item
  Our goal is to train a model that can also predict well on new, unseen
  data
\item
  We simulate the situation of unseen data by splitting our data into
  one part for model selection and another part for model evaluation
\item
  Model evaluation is usually done via Cross Validation, but classic
  Cross Validation does not take into account the time dependency in our
  data
\end{itemize}
\end{frame}

\begin{frame}{Forward selection}
\protect\hypertarget{forward-selection}{}
\begin{center}\includegraphics[width=0.75\linewidth]{../thesis/validation-schemes} \end{center}

\begin{center}\includegraphics[width=0.75\linewidth]{../thesis/blocked-cv} \end{center}
\end{frame}

\begin{frame}{Forward selection}
\protect\hypertarget{forward-selection-1}{}
\begin{itemize}
\tightlist
\item
  We compute a \(\lambda\)-vector for the complete training set
\item
  For each fold we fit a model with this \(\lambda\)-vector
\item
  We compute the prediction error for the cv-test set of each fold
\item
  Choose \(\lambda_{\min}\), \(\lambda\) that minimizes average MSE over
  all folds
\item
  Fit model on complete selection data with \(\lambda_{\min}\) and
  compute MSE on evaluation data
\end{itemize}
\end{frame}

\begin{frame}{lasso settings}
\protect\hypertarget{lasso-settings}{}
\begin{itemize}
\tightlist
\item
  lasso
\item
  lasso with standardized features
\item
  lasso with de-seasonalized SST
\item
  lasso with differentiated SST
\item
  lasso on clusters
\end{itemize}
\end{frame}

\begin{frame}{lasso results TODO}
\protect\hypertarget{lasso-results-todo}{}
\begin{itemize}
\tightlist
\item
  Show only best model results
\item
  lasso with standardized features
\item
  show MSE in plots
\item
  show predictions in plots
\item
  show predictions
\item
  show coefficients
\item
  display table
\end{itemize}
\end{frame}

\begin{frame}{Lasso results}
\protect\hypertarget{lasso-results}{}
\end{frame}

\begin{frame}{MSE in each fold}
\protect\hypertarget{mse-in-each-fold}{}
\begin{center}\includegraphics[width=0.75\linewidth]{ma-presentation_files/figure-beamer/unnamed-chunk-16-1} \end{center}
\end{frame}

\begin{frame}{Predictions on test set, for each Fold}
\protect\hypertarget{predictions-on-test-set-for-each-fold}{}
\begin{figure}

{\centering \includegraphics[width=0.75\linewidth]{ma-presentation_files/figure-beamer/pred-plot-fold-lasso-stand-1} 

}

\caption{Precipitation prediction and target values in the test set in each fold. Predictions in red and target values in black.}\label{fig:pred-plot-fold-lasso-stand}
\end{figure}
\end{frame}

\hypertarget{predictions-on-external-test-set}{%
\section{Predictions on External Test
Set}\label{predictions-on-external-test-set}}

\begin{frame}[fragile]{Predictions on External Test Set}
\begin{figure}

{\centering \includegraphics[width=0.75\linewidth]{ma-presentation_files/figure-beamer/pred-plot-full-lasso-stand-1} 

}

\caption{Precipitation prediction and target values in the validation set. Predictions in red and target values in black. The model was fitted on the full CV data with the lambda value that minimised the average MSE}\label{fig:pred-plot-full-lasso-stand}
\end{figure}

\begin{verbatim}
## [1] 1214.49
\end{verbatim}
\end{frame}

\begin{frame}{SST Regions chosen by the lasso}
\protect\hypertarget{sst-regions-chosen-by-the-lasso}{}
\begin{figure}

{\centering \includegraphics[width=0.75\linewidth]{ma-presentation_files/figure-beamer/coef-plot-full-lasso-stand-1} 

}

\caption{Coefficient plot of the full lasso model.}\label{fig:coef-plot-full-lasso-stand}
\end{figure}
\end{frame}

\hypertarget{lasso-results-all-models}{%
\section{Lasso results all models}\label{lasso-results-all-models}}

\begin{frame}{Lasso results all models}
\begin{longtable}[]{@{}lrr@{}}
\toprule()
& MSE & Lambda \\
\midrule()
\endhead
Standardized & 1214.49 & 3.52 \\
Original & 1314.93 & 3.52 \\
Differentiated & 1361.82 & 2.21 \\
De-seasonalized & 1809.45 & 1.75 \\
\bottomrule()
\end{longtable}
\end{frame}

\begin{frame}{Lasso on clusters results}
\protect\hypertarget{lasso-on-clusters-results}{}
\begin{center}\includegraphics[width=0.5\linewidth]{ma-presentation_files/figure-beamer/cluster-map2-1} \end{center}
\end{frame}

\begin{frame}{MSE clustered results}
\protect\hypertarget{mse-clustered-results}{}
\begin{center}\includegraphics[width=0.75\linewidth]{ma-presentation_files/figure-beamer/unnamed-chunk-20-1} \end{center}
\end{frame}

\begin{frame}{Summary lasso results}
\protect\hypertarget{summary-lasso-results}{}
\begin{itemize}
\tightlist
\item
  We compared different settings for the lasso
\item
  Lasso, lasso with standardization, de-seasonalizing, differentiation
  and on clusters
\item
  Lasso with standardized SST worked best
\item
  Can predict general seasonality, but still fails to predict peaks in
  precipitation
\item
  Clustering the CAB improves only on one cluster, but on this cluster
  peaks can be predicted better than in the unclustered model
\item
  Lasso chooses single ``points'' and not whole areas
\item
  The points chosen as coefficients differ in the models, and can be
  very far away from the CAB
\end{itemize}
\end{frame}

\begin{frame}{Discussion Validation approach (maybe discuss this at the
end)}
\protect\hypertarget{discussion-validation-approach-maybe-discuss-this-at-the-end}{}
\begin{itemize}
\tightlist
\item
  For the CAB we can not predict large values in the hold-out set, on
  cluster 2 it works a little better
\item
  Possible explanations:
\item
  Our validation approach works better when train and test set are
  similar in terms of seasonality and trend
\item
  When train and test differ, predictions might not work so well (test
  of stationarity in folds)
\item
  Differentiating and de-seasonalizing could not solve this problem
\item
  Predictions work better when the precipitation remains fairly stable
  over time, see Cluster 2
\item
  Final model uses complete model selection data, possibly some of that
  information is not useful anymore if it's toocfar away from hold-out
  time frame
\item
  Our validation approach is a trade-off between efficient use of data
  and acknowledging the time-dependence in the data
\item
  Possible challenges are non-stationarity and spurious correlations
\item
  Lasso chooses single points, if the predictive area changes, its
  coefficient is 0
\end{itemize}
\end{frame}

\hypertarget{the-fused-lasso}{%
\section{The fused lasso}\label{the-fused-lasso}}

\begin{frame}{Definition of the fused lasso}
\protect\hypertarget{definition-of-the-fused-lasso}{}
\begin{itemize}
\tightlist
\item
  Fused lasso, ``fuses'' predictors together
\item
  It penalizes the difference of close predictores
\item
  Therefore close predictors should be similar
\end{itemize}

\begin{equation}
\min_{\beta} 1/2 \sum_{i=1}^n(y_i - x_i^T\beta_i)^2 + \lambda \sum_{i,j \in E} |\beta_i - \beta_j| + \gamma \cdot \lambda \sum_{i=1}^p|\beta_i|,
(\#eq:fused-lasso)
\end{equation}

\begin{itemize}
\tightlist
\item
  with \(x_i\) being the ith row of the predictor matrix and E is the
  edge set of an underlying graph.
\item
  The third term \(\gamma \cdot \lambda \sum_{i=1}^p|\beta_i|\),
  controls the sparsity of the coefficients.
\item
  \(\gamma=0\) leads to complete fusion of the coefficients (no
  sparsity) and \(\gamma\) \textgreater{} 0 introduces sparsity to the
  solution, with higher values placing more priority on sparsity.
\end{itemize}
\end{frame}

\begin{frame}{Fused lasso optimization}
\protect\hypertarget{fused-lasso-optimization}{}
\begin{itemize}
\tightlist
\item
  Lets consider the problem in the notation of the generalized lasso
  problem
\end{itemize}

\begin{equation}
\hat{\beta}=\underset{\beta \in \mathbb{R}^{p}}{\arg \min } \frac{1}{2}\|y-X \beta\|_{2}^{2}+\lambda\|D \beta\|_{1},
(\#eq:gen-lasso)
\end{equation}

\begin{itemize}
\tightlist
\item
  where \(y \in \mathbb{R}^n\) is the vector of the outcome,
  \(X \in \mathbb{R}^{n \times p}\) a predictor matrix,
  \(D \in \mathbb{R}^{m \times p}\) denotes a penalty matrix, and
  \(\lambda \geq 0\) is a regularization parameter.
\item
  The dual path algorithm solves not the primal but the dual solution of
  the problem and computes the solution for a whole path instead of
  single values of \(\lambda\).
\end{itemize}
\end{frame}

\begin{frame}{}
\protect\hypertarget{section}{}
\begin{itemize}
\tightlist
\item
  Let's consider the case when \(X=I\) and \(rank(X)=p\) (this is called
  the ``signal approximator'' case), the dual problem of
  @ref(eq:gen-lasso) is then:
\end{itemize}

\begin{equation}
\hat{u} \in \underset{\substack{u \in \mathbb{R}^{\omega}}}{\arg \min } \frac{1}{2}\left\|y-D^{T} u\right\|_{\frac{2}{2}} \text { subject to }\|u\|_{\infty} \leq \lambda.
(\#eq:dual)
\end{equation}

\begin{itemize}
\tightlist
\item
  The primal and dual solutions, \(\hat{\beta}\) and \(\hat{u}\) are
  related by:
\end{itemize}

\begin{equation}
\hat{\beta}=y-D^{T} \hat{u} .
(\#eq:dual-relate)
\end{equation}

\begin{itemize}
\tightlist
\item
  For general X and D with exploitable structure (as in our case),
  specialized implementations exist
\end{itemize}
\end{frame}

\begin{frame}[fragile]{Graph structure}
\protect\hypertarget{graph-structure}{}
\begin{itemize}
\tightlist
\item
  We can use a graph as input in the fusedlasso function
\item
  We created a grid and deleted all nodes that were land regions
\item
  This induced subgraphs
\end{itemize}

\begin{figure}

{\centering \includegraphics[width=0.75\linewidth]{ma-presentation_files/figure-beamer/graph-plot2-1} 

}

\caption{Graph of the SST and land areas used in fused lasso}\label{fig:graph-plot2}
\end{figure}

\begin{verbatim}
## NULL
\end{verbatim}
\end{frame}

\begin{frame}{Graph structure and implications}
\protect\hypertarget{graph-structure-and-implications}{}
\begin{itemize}
\tightlist
\item
  Results showed that removing the sub-graphs improved performance,
  although some of the regions were included in the final lasso models
\item
  If we don't remove the clusters and also add sparsity (i.e
  \(\gamma > 0\)) the clusters dominate the results even more
\item
  Possible explanations: Sub-graphs are less penalized, because they
  have fewer edges.
\item
  Removing the clusters improved results more than f.e standardization
\end{itemize}
\end{frame}

\begin{frame}{Fused lasso settings}
\protect\hypertarget{fused-lasso-settings}{}
\begin{itemize}
\tightlist
\item
  The considered fused lasso settings are: Fused lasso with clusters,
  fused lasso without clusters, fused lasso without clusters and
  sparsity (gamma: 0.01, 0.05, 0.1)
\item
  Fused lasso without clusters and no sparsity showed best results
\end{itemize}
\end{frame}

\begin{frame}{Fused lasso results, clusters removed}
\protect\hypertarget{fused-lasso-results-clusters-removed}{}
\begin{center}\includegraphics[width=0.75\linewidth]{err-line-plot-noclust} \end{center}
\end{frame}

\begin{frame}{Fused lasso results, clusters removed}
\protect\hypertarget{fused-lasso-results-clusters-removed-1}{}
\begin{center}\includegraphics[width=0.75\linewidth]{../results/CV-fused/noclust-large-fused-5k/best-lambda-res-noclust} \end{center}
\end{frame}

\begin{frame}{Prediction plots}
\protect\hypertarget{prediction-plots}{}
\begin{figure}

{\centering \includegraphics[width=0.75\linewidth]{ma-presentation_files/figure-beamer/pred-fold-fused-og-1} 

}

\caption{Precipitation prediction and target values in the test set in each fold. Predictions in red and target values in black.}\label{fig:pred-fold-fused-og}
\end{figure}
\end{frame}

\begin{frame}{Prediction plot}
\protect\hypertarget{prediction-plot}{}
The predictions inside the folds are very similar to lasso without
standardization, the same holds for the predictions from the full model,
but the MSE improves here.
\end{frame}

\begin{frame}[fragile]{Full predictions}
\protect\hypertarget{full-predictions}{}
\begin{figure}

{\centering \includegraphics[width=0.75\linewidth]{ma-presentation_files/figure-beamer/pred-plot-full-fused-og-1} 

}

\caption{Precipitation prediction and target values in the validation set. Predictions in red and target values in black. The model was fitted on the full CV data with the lambda value that minimised the average MSE}\label{fig:pred-plot-full-fused-og}
\end{figure}

\begin{verbatim}
## [1] 1070.04
\end{verbatim}
\end{frame}

\begin{frame}{Coefficient plot}
\protect\hypertarget{coefficient-plot}{}
\begin{center}\includegraphics[width=0.75\linewidth]{../results/CV-fused/noclust-large-fused-5k/coef-plots/coef-plot-full} \end{center}
\end{frame}

\begin{frame}{Coefficient plot, highest absolute values only}
\protect\hypertarget{coefficient-plot-highest-absolute-values-only}{}
\begin{center}\includegraphics[width=0.75\linewidth]{../results/CV-fused/noclust-large-fused-5k/coef-plots/coef-plot-drop-out-full1} \end{center}
\end{frame}

\begin{frame}{Coefficient plot, lowest absolute values only}
\protect\hypertarget{coefficient-plot-lowest-absolute-values-only}{}
\begin{center}\includegraphics[width=0.75\linewidth]{../results/CV-fused/noclust-large-fused-5k/coef-plots/coef-plot-drop-out-full2} \end{center}
\end{frame}

\begin{frame}{Fused lasso results MSE}
\protect\hypertarget{fused-lasso-results-mse}{}
\begin{longtable}[]{@{}lrr@{}}
\toprule()
& MSE & Lambda \\
\midrule()
\endhead
No sub-graphs & 1070.04 & 14.28 \\
With sub-graphs & 1131.71 & 18.55 \\
No sub-graphs, gamma 0.05 & 1836.63 & 2544.23 \\
No sub-graphs, gamma 0.1 & 1840.59 & 1586.29 \\
\bottomrule()
\end{longtable}
\end{frame}

\begin{frame}{Fused lasso results summary}
\protect\hypertarget{fused-lasso-results-summary}{}
\begin{itemize}
\tightlist
\item
  We compared different settings for the fused lasso, removing the
  sub-graphs and introducing no sparsity gave the best results
\item
  Removing the sub-graphs removed some of the optimization problems, but
  nodes with less edges are still less penalized
\item
  Implementing a validation strategy was more complex than for the lasso
\item
  We smoothed the error-lines in each fold over a common region to
  compute \(\lambda_{\min}\)
\item
  The coefficient plots reveal predictive areas with high negative
  values in the Baltic Sea and high positive values north east of
  Canada.
\item
  Since no sparsity is used, all areas obtain non-zero coefficient
  values
\end{itemize}
\end{frame}

\begin{frame}{Fused lasso discussion}
\protect\hypertarget{fused-lasso-discussion}{}
\begin{itemize}
\tightlist
\item
  Computing the solution path is computationally expensive
\item
  The graph structure is highly influential and cost will scale with
  number of edges
\item
  While the best fused lasso approach performed best overall, it still
  is not able to predict high precipitation values
\item
  Possible improvements on optimization path: creating weighted graph
  (increases number of edges and cost), narrowing down the SST
  ``window'' (f.e as in Ciemer et al. (2020))
\item
  Possible improvements on feature engineering: as for the lasso,
  differentiating, de-seasonalizing
\end{itemize}
\end{frame}

\hypertarget{discussion-conclusion}{%
\section{Discussion \& Conclusion}\label{discussion-conclusion}}

\begin{frame}{Discussion}
\protect\hypertarget{discussion}{}
\begin{itemize}
\tightlist
\item
  Our results suggest that precipitation can to some extend be predicted
  from SST directly.
\item
  The overall predictability of precipitation in the CAB differed
  between model selection and model evaluation phase.
\item
  For one part this might be due to the difference in the regions in the
  CAB, since clustering improved the results for one specific cluster.
\item
  Another explanation could be that our model selection approach was not
  optimal in its use of the data.
\item
  We might have been to restrictive in exploiting the data or used data
  that became less relevant over time.
\end{itemize}
\end{frame}

\begin{frame}{Discussion}
\protect\hypertarget{discussion-1}{}
\begin{itemize}
\tightlist
\item
  Possible other approaches:
\item
  Allow for larger folds (introduces overlapping folds), or for crossing
  of train and test in time (past is predicted with future values)
\item
  Fit the full model with less data and discarding data that is far away
  from the hold-out validation set.
\item
  The results of the fused lasso will depend a lot on the graph
  structure, sub-graphs do not represent the real situation well
\item
  Creating a weighted graph or narrowing the SST ``window'' may improve
  performance.
\item
  Also, applying the fused lasso only on the best performing cluster
  from the lasso may yield better results.
\end{itemize}
\end{frame}

\begin{frame}{Conclusion}
\protect\hypertarget{conclusion}{}
\begin{itemize}
\tightlist
\item
  In a descriptive analysis we found temporal and spatial patterns in
  the correlation of rain in the CAB and SST
\item
  The cluster analysis revealed 5 almost completely spatially coherent
  clusters in the CAB
\item
  Standardizing the features yielded the best results for the lasso
\item
  The lasso can predict the precipitation on the model selection test
  sets a lot better than in the hold-out test set
\item
  On the hold-out data the lasso fails to predict the peaks in
  precipitation
\end{itemize}
\end{frame}

\begin{frame}{Conclusion}
\protect\hypertarget{conclusion-1}{}
\begin{itemize}
\tightlist
\item
  We applied the fused lasso to our problem and implemented a model
  evaluation approach
\item
  The fused lasso improves predictive power compared to the lasso when
  the sub-graphs are removed
\item
  The fused lasso is still is not able to predict high values in
  precipitation well
\item
  We could further improve the clustering method by taking into account
  spatial dependencies.
\item
  The fused lasso could be improved by using other model selection
  approaches or increasing the complexity of the graph structure.
\end{itemize}
\end{frame}

\hypertarget{appendix}{%
\section{Appendix}\label{appendix}}

\hypertarget{lasso-on-original-sst}{%
\section{Lasso on original SST}\label{lasso-on-original-sst}}

\begin{frame}{MSE in each fold (Lasso on original SST)}
\protect\hypertarget{mse-in-each-fold-lasso-on-original-sst}{}
\begin{center}\includegraphics[width=0.75\linewidth]{ma-presentation_files/figure-beamer/unnamed-chunk-34-1} \end{center}
\end{frame}

\begin{frame}{Predictions on test set, for each Fold (Lasso original
SST)}
\protect\hypertarget{predictions-on-test-set-for-each-fold-lasso-original-sst}{}
\begin{figure}

{\centering \includegraphics[width=0.75\linewidth]{ma-presentation_files/figure-beamer/pred-plot-fold-lasso-og-1} 

}

\caption{Precipitation prediction and target values in the test set in each fold. Predictions in red and target values in black.}\label{fig:pred-plot-fold-lasso-og}
\end{figure}
\end{frame}

\begin{frame}[fragile]{Predictions on External Test Set (Lasso on
original SST)}
\protect\hypertarget{predictions-on-external-test-set-lasso-on-original-sst}{}
\begin{figure}

{\centering \includegraphics[width=0.75\linewidth]{ma-presentation_files/figure-beamer/pred-plot-full-lasso-og-1} 

}

\caption{Precipitation prediction and target values in the validation set. Predictions in red and target values in black. The model was fitted on the full CV data with the lambda value that minimised the average MSE}\label{fig:pred-plot-full-lasso-og}
\end{figure}

\begin{verbatim}
## [1] 1314.93
\end{verbatim}
\end{frame}

\begin{frame}{SST Regions chosen by the lasso (Lasso on original SST)}
\protect\hypertarget{sst-regions-chosen-by-the-lasso-lasso-on-original-sst}{}
\begin{figure}

{\centering \includegraphics[width=0.75\linewidth]{ma-presentation_files/figure-beamer/coef-plot-full-lasso-og-1} 

}

\caption{Coefficient plot of the full lasso model.}\label{fig:coef-plot-full-lasso-og}
\end{figure}
\end{frame}

\hypertarget{lasso-on-differentiated-sst}{%
\section{Lasso on differentiated
SST}\label{lasso-on-differentiated-sst}}

\begin{frame}{MSE in each fold (Lasso differentiated)}
\protect\hypertarget{mse-in-each-fold-lasso-differentiated}{}
\begin{center}\includegraphics[width=0.75\linewidth]{ma-presentation_files/figure-beamer/unnamed-chunk-38-1} \end{center}
\end{frame}

\begin{frame}{Predictions on test set, for each Fold (Lasso
differentiated)}
\protect\hypertarget{predictions-on-test-set-for-each-fold-lasso-differentiated}{}
\begin{figure}

{\centering \includegraphics[width=0.75\linewidth]{ma-presentation_files/figure-beamer/pred-plot-fold-lasso-diff-1} 

}

\caption{Precipitation prediction and target values in the test set in each fold. Predictions in red and target values in black.}\label{fig:pred-plot-fold-lasso-diff}
\end{figure}
\end{frame}

\begin{frame}[fragile]{Predictions on External Test Set (Lasso
differentiated)}
\protect\hypertarget{predictions-on-external-test-set-lasso-differentiated}{}
\begin{figure}

{\centering \includegraphics[width=0.75\linewidth]{ma-presentation_files/figure-beamer/pred-plot-full-lasso-diff-1} 

}

\caption{Precipitation prediction and target values in the validation set. Predictions in red and target values in black. The model was fitted on the full CV data with the lambda value that minimised the average MSE}\label{fig:pred-plot-full-lasso-diff}
\end{figure}

\begin{verbatim}
## [1] 1361.82
\end{verbatim}
\end{frame}

\begin{frame}{SST Regions chosen by the lasso (Lasso differentiated)}
\protect\hypertarget{sst-regions-chosen-by-the-lasso-lasso-differentiated}{}
\begin{figure}

{\centering \includegraphics[width=0.75\linewidth]{ma-presentation_files/figure-beamer/coef-plot-full-lasso-diff-1} 

}

\caption{Coefficient plot of the full lasso model.}\label{fig:coef-plot-full-lasso-diff}
\end{figure}
\end{frame}

\begin{frame}{Fused evaluation (maybe explain this when showing
resutls)}
\protect\hypertarget{fused-evaluation-maybe-explain-this-when-showing-resutls}{}
\begin{itemize}
\tightlist
\item
  Generally same setting as for lasso, 5 folds with train and test,
  choose \(\lambda_{\min}\), refit with \(\lambda_{\min}\), get MSE on
  hold-out test set.
\item
  But for the fused lasso we can not define the \(\lambda\) vector
  beforehand.
\item
  \(\lambda\)-path is found by dual path algorithm and the range of the
  paths can vary a lot!
\item
  So to find \(\lambda_{\min}\) we search over the common range of all
  folds and interpolate to lines
\item
  \(\lambda_{\min}\) is then the \(\lambda\) that minimize MSE over all
  \(\lambda\) of that common range
\end{itemize}
\end{frame}

\hypertarget{lasso-on-de-seasonalized-sst}{%
\section{Lasso on de-seasonalized
SST}\label{lasso-on-de-seasonalized-sst}}

\begin{frame}{MSE in each fold (Lasso de-seasonalized)}
\protect\hypertarget{mse-in-each-fold-lasso-de-seasonalized}{}
\begin{center}\includegraphics[width=0.75\linewidth]{ma-presentation_files/figure-beamer/unnamed-chunk-42-1} \end{center}
\end{frame}

\begin{frame}{Predictions on test set, for each Fold (Lasso
de-seasonalized)}
\protect\hypertarget{predictions-on-test-set-for-each-fold-lasso-de-seasonalized}{}
\begin{figure}

{\centering \includegraphics[width=0.75\linewidth]{ma-presentation_files/figure-beamer/pred-plot-fold-lasso-des-1} 

}

\caption{Precipitation prediction and target values in the test set in each fold. Predictions in red and target values in black.}\label{fig:pred-plot-fold-lasso-des}
\end{figure}
\end{frame}

\begin{frame}[fragile]{Predictions on External Test Set (Lasso
differentiated)}
\protect\hypertarget{predictions-on-external-test-set-lasso-differentiated-1}{}
\begin{figure}

{\centering \includegraphics[width=0.75\linewidth]{ma-presentation_files/figure-beamer/pred-plot-full-lasso-des-1} 

}

\caption{Precipitation prediction and target values in the validation set. Predictions in red and target values in black. The model was fitted on the full CV data with the lambda value that minimised the average MSE}\label{fig:pred-plot-full-lasso-des}
\end{figure}

\begin{verbatim}
## [1] 1809.45
\end{verbatim}
\end{frame}

\begin{frame}{SST Regions chosen by the lasso (Lasso de-seasonalized)}
\protect\hypertarget{sst-regions-chosen-by-the-lasso-lasso-de-seasonalized}{}
\begin{figure}

{\centering \includegraphics[width=0.75\linewidth]{ma-presentation_files/figure-beamer/coef-plot-full-lasso-des-1} 

}

\caption{Coefficient plot of the full lasso model.}\label{fig:coef-plot-full-lasso-des}
\end{figure}
\end{frame}

\hypertarget{fused-lasso-with-sub-graphs}{%
\section{Fused lasso with
sub-graphs}\label{fused-lasso-with-sub-graphs}}

\begin{frame}{Error lines (Fused lasso with sub-graphs)}
\protect\hypertarget{error-lines-fused-lasso-with-sub-graphs}{}
\begin{center}\includegraphics[width=0.45\linewidth]{../results/CV-fused/large-fused-5k/best-lambda-res} \end{center}
\end{frame}

\begin{frame}{Prediction plots for each fold (Fused lasso with
sub-graphs)}
\protect\hypertarget{prediction-plots-for-each-fold-fused-lasso-with-sub-graphs}{}
\begin{figure}

{\centering \includegraphics[width=0.75\linewidth]{ma-presentation_files/figure-beamer/pred-fold-fused-sub-1} 

}

\caption{Precipitation prediction and target values in the test set in each fold. Predictions in red and target values in black.}\label{fig:pred-fold-fused-sub}
\end{figure}

The predictions inside the folds are very similar to lasso without
standardization (see @ref(fig:pred-plot-fold-lasso-og)), the same holds
for the predictions from the full model, but the MSE improves here
(@\ref(fig:pred-plot-full-fused-og).
\end{frame}

\begin{frame}[fragile]{Predictions on hold-out set (Fused lasso with
sub-graphs)}
\protect\hypertarget{predictions-on-hold-out-set-fused-lasso-with-sub-graphs}{}
\begin{figure}

{\centering \includegraphics[width=0.75\linewidth]{ma-presentation_files/figure-beamer/pred-plot-full-fused-sub-1} 

}

\caption{Precipitation prediction and target values in the validation set. Predictions in red and target values in black. The model was fitted on the full CV data with the lambda value that minimised the average MSE}\label{fig:pred-plot-full-fused-sub}
\end{figure}

\begin{verbatim}
## [1] 1131.71
\end{verbatim}
\end{frame}

\begin{frame}{Coefficient plot of the final model (Fused lasso with
sub-graphs)}
\protect\hypertarget{coefficient-plot-of-the-final-model-fused-lasso-with-sub-graphs}{}
\begin{center}\includegraphics[width=0.75\linewidth]{../results/CV-fused/large-fused-5k/coef-plots/coef-plot-full} \end{center}

\begin{center}\includegraphics[width=0.75\linewidth]{../results/CV-fused/large-fused-5k/coef-plots/coef-plot-drop-out-full1} \end{center}

\begin{center}\includegraphics[width=0.75\linewidth]{../results/CV-fused/large-fused-5k/coef-plots/coef-plot-drop-out-full2} \end{center}
\end{frame}

\hypertarget{fused-lasso-without-sub-graphs-gamma-0.1}{%
\section{Fused lasso without sub-graphs, gamma
0.1}\label{fused-lasso-without-sub-graphs-gamma-0.1}}

\begin{frame}{Error lines (Fused lasso without sub-graphs, gamma 0.1)}
\protect\hypertarget{error-lines-fused-lasso-without-sub-graphs-gamma-0.1}{}
\begin{center}\includegraphics[width=0.45\linewidth]{../results/CV-fused/noclust-large-fused-5k-gamma-01/best-lambda-res} \end{center}
\end{frame}

\begin{frame}{Prediction plots for each fold (Fused lasso without
sub-graphs, gamma 0.1)}
\protect\hypertarget{prediction-plots-for-each-fold-fused-lasso-without-sub-graphs-gamma-0.1}{}
\begin{figure}

{\centering \includegraphics[width=0.75\linewidth]{ma-presentation_files/figure-beamer/pred-fold-fused-nolcust-01-1} 

}

\caption{Precipitation prediction and target values in the test set in each fold. Predictions in red and target values in black.}\label{fig:pred-fold-fused-nolcust-01}
\end{figure}

The predictions inside the folds are very similar to lasso without
standardization (see @ref(fig:pred-plot-fold-lasso-og)), the same holds
for the predictions from the full model, but the MSE improves here
(@\ref(fig:pred-plot-full-fused-og).
\end{frame}

\begin{frame}[fragile]{Predictions on hold-out set (Fused lasso without
sub-graphs, gamma 0.1)}
\protect\hypertarget{predictions-on-hold-out-set-fused-lasso-without-sub-graphs-gamma-0.1}{}
\begin{figure}

{\centering \includegraphics[width=0.75\linewidth]{ma-presentation_files/figure-beamer/pred-plot-full-fused-noclust-01-1} 

}

\caption{Precipitation prediction and target values in the validation set. Predictions in red and target values in black. The model was fitted on the full CV data with the lambda value that minimised the average MSE}\label{fig:pred-plot-full-fused-noclust-01}
\end{figure}

\begin{verbatim}
## [1] 1840.59
\end{verbatim}
\end{frame}

\begin{frame}{Coefficient plot of the final model (Fused lasso without
sub-graphs, gamma 0.1)}
\protect\hypertarget{coefficient-plot-of-the-final-model-fused-lasso-without-sub-graphs-gamma-0.1}{}
\begin{center}\includegraphics[width=0.75\linewidth]{../results/CV-fused/noclust-large-fused-5k-gamma-01/coef-plots/coef-plot-full} \end{center}

\begin{center}\includegraphics[width=0.75\linewidth]{../results/CV-fused/noclust-large-fused-5k-gamma-01/coef-plots/coef-plot-drop-out-full1} \end{center}

\begin{center}\includegraphics[width=0.75\linewidth]{../results/CV-fused/noclust-large-fused-5k-gamma-01/coef-plots/coef-plot-drop-out-full2} \end{center}
\end{frame}

\hypertarget{fused-lasso-without-sub-graphs-gamma-0.05}{%
\section{Fused lasso without sub-graphs, gamma
0.05}\label{fused-lasso-without-sub-graphs-gamma-0.05}}

\begin{frame}{Error lines (Fused lasso without sub-graphs, gamma 0.05)}
\protect\hypertarget{error-lines-fused-lasso-without-sub-graphs-gamma-0.05}{}
\begin{center}\includegraphics[width=0.45\linewidth]{../results/CV-fused/noclust-large-fused-5k-gamma-005/best-lambda-res} \end{center}
\end{frame}

\begin{frame}{Prediction plots for each fold (Fused lasso without
sub-graphs, gamma 0.05)}
\protect\hypertarget{prediction-plots-for-each-fold-fused-lasso-without-sub-graphs-gamma-0.05}{}
\begin{figure}

{\centering \includegraphics[width=0.75\linewidth]{ma-presentation_files/figure-beamer/pred-fold-fused-noclust-005-1} 

}

\caption{Precipitation prediction and target values in the test set in each fold. Predictions in red and target values in black.}\label{fig:pred-fold-fused-noclust-005}
\end{figure}

The predictions inside the folds are very similar to lasso without
standardization (see @ref(fig:pred-plot-fold-lasso-og)), the same holds
for the predictions from the full model, but the MSE improves here
(@\ref(fig:pred-plot-full-fused-og).
\end{frame}

\begin{frame}[fragile]{Predictions on hold-out set (Fused lasso without
sub-graphs, gamma 0.1)}
\protect\hypertarget{predictions-on-hold-out-set-fused-lasso-without-sub-graphs-gamma-0.1-1}{}
\begin{figure}

{\centering \includegraphics[width=0.75\linewidth]{ma-presentation_files/figure-beamer/pred-plot-full-fused-noclust-05-1} 

}

\caption{Precipitation prediction and target values in the validation set. Predictions in red and target values in black. The model was fitted on the full CV data with the lambda value that minimised the average MSE}\label{fig:pred-plot-full-fused-noclust-05}
\end{figure}

\begin{verbatim}
## [1] 1836.63
\end{verbatim}
\end{frame}

\begin{frame}{Coefficient plot of the final model (Fused lasso without
sub-graphs, gamma 0.05)}
\protect\hypertarget{coefficient-plot-of-the-final-model-fused-lasso-without-sub-graphs-gamma-0.05}{}
\begin{center}\includegraphics[width=0.75\linewidth]{../results/CV-fused/noclust-large-fused-5k-gamma-005/coef-plots/coef-plot-full} \end{center}

\begin{center}\includegraphics[width=0.75\linewidth]{../results/CV-fused/noclust-large-fused-5k-gamma-005/coef-plots/coef-plot-drop-out-full1} \end{center}

\begin{center}\includegraphics[width=0.75\linewidth]{../results/CV-fused/noclust-large-fused-5k-gamma-005/coef-plots/coef-plot-drop-out-full2} \end{center}
\end{frame}

\hypertarget{fused-evaluation-maybe-explain-this-when-showing-resutls-1}{%
\section{Fused evaluation (maybe explain this when showing
resutls)}\label{fused-evaluation-maybe-explain-this-when-showing-resutls-1}}

\begin{frame}{Fused evaluation (maybe explain this when showing
resutls)}
\begin{itemize}
\tightlist
\item
  Generally same setting as for lasso, 5 folds with train and test,
  choose \(\lambda_{\min}\), refit with \(\lambda_{\min}\), get MSE on
  hold-out test set.
\item
  But for the fused lasso we can not define the \(\lambda\) vector
  beforehand.
\item
  \(\lambda\)-path is found by dual path algorithm and the range of the
  paths can vary a lot!
\item
  So to find \(\lambda_{\min}\) we search over the common range of all
  folds and interpolate to lines
\item
  \(\lambda_{\min}\) is then the \(\lambda\) that minimize MSE over all
  \(\lambda\) of that common range
\end{itemize}

\hypertarget{refs}{}
\begin{CSLReferences}{1}{0}
\leavevmode\vadjust pre{\hypertarget{ref-ciemer2020early}{}}%
Ciemer, Catrin, Lars Rehm, Juergen Kurths, Reik V Donner, Ricarda
Winkelmann, and Niklas Boers. 2020. {``An Early-Warning Indicator for
Amazon Droughts Exclusively Based on Tropical Atlantic Sea Surface
Temperatures.''} \emph{Environmental Research Letters} 15 (9): 094087.

\end{CSLReferences}
\end{frame}

\end{document}
